\subsection{Netduino Plus 2 Initialization \& Pin Connection Architecture}

The \texttt{netduinoplus2\_init()} (Code~\ref{code:netduinoplus2_init}) function, is responsible for the complete system 
assembly, following QEMU's device instantiation hierarchy. The first thing this 
function does is instantiate the STM32F405 SoC using the \texttt{qdev\_new()} function with the parameter \texttt{TYPE\_STM32F405\_SOC} 
and realize it with the \texttt{sysbus\_realize\_} \texttt{and\_unref()} function. 
This last function activates the SoC's internal \texttt{stm32f405\_} \texttt{soc\_realize()} 
function, which instantiates the three I$^2$C controllers and populates the 
\texttt{s->i2c[0].bus} pointer with the I2C1 bus reference, making it available 
for connecting external devices. 

Once the SoC is implemented, the LIS3DH accelerometer is created again using 
the \texttt{qdev\_new()} function, but with the \texttt{TYPE\_LIS3DH}. It is configured with its 7-bit 
I2C address \texttt{0x18} (pin SA0 connected to ground) and connected to I2C1 using the 
\texttt{i2c\_slave\_realize\_and\_unref()} function.


The LIS3DH's INT1 interrupt output is routed to the STM32's EXTI 1 line using 
\texttt{qdev\_connect\_gpio\_out()}, allowing the firmware to receive data-ready 
notifications. Finally, \texttt{armv7m\_load\_kernel()} loads the user-supplied 
firmware binary into the emulated flash memory, resulting in a fully functional 
Netduino Plus 2 emulation environment.

\newpage

\begin{lstlisting}[language=C, label={code:netduinoplus2_init}, caption={Netduino Plus 2 machine initialization}]
static void netduinoplus2_init(MachineState *machine)
{
    DeviceState *dev;
    Clock *sysclk;

    /* This clock doesn't need migration because it is fixed-frequency */
    sysclk = clock_new(OBJECT(machine), "SYSCLK");
    clock_set_hz(sysclk, SYSCLK_FRQ);

    /* Instantiate STM32F405 SoC */
    dev = qdev_new(TYPE_STM32F405_SOC);
    object_property_add_child(OBJECT(machine), "soc", OBJECT(dev));
    qdev_connect_clock_in(dev, "sysclk", sysclk);
    sysbus_realize_and_unref(SYS_BUS_DEVICE(dev), &error_fatal);

    /* Initialize LIS3DH on I2C1 */
    STM32F405State *s = STM32F405_SOC(dev);    
    DeviceState *lis3dh = qdev_new(TYPE_LIS3DH); // Creating LIS3DH accelerometer
    
    i2c_slave_set_address(I2C_SLAVE(lis3dh), 0x18); // Setting I2C address to 0x18    
    i2c_slave_realize_and_unref(I2C_SLAVE(lis3dh), 
                    s->i2c[0].bus, &error_fatal);    // Connecting to I2C1 bus

    qdev_connect_gpio_out(lis3dh, 0, qdev_get_gpio_in(DEVICE(&s->exti), 1));

    /* Load firmware into flash memory */
    armv7m_load_kernel(ARM_CPU(first_cpu),
                       machine->kernel_filename,
                       0, FLASH_SIZE);
}
\end{lstlisting}


\subsection{QEMU Device Hierarchy}

When a Netduino Plus 2 machine instance is created, a hierarchical device tree is also created where the LIS3DH exists as an I$^2$C slave device connected to the I2C1 bus of the STM32F405 SoC. Figure ~\ref{fig:qemu_hierarchy} illustrates the complete device hierarchy:

\begin{figure}[H]
    \centering
    \begin{verbatim}
    netduinoplus2 (Machine)
    \---- STM32F405_SOC
        |---- I2C[0] (I2C Bus)
        |   \---- lis3dh (I2C Slave @ 0x18)
        |       |---- int1 (qemu_irq out)  ----\
        |       \---- int2 (qemu_irq out)      |
        |                                      |
        \---- exti (STM32F4XX_EXTI)            |
            |---- gpio_in[0] ← (EXTI0)         |
            |---- gpio_in[1] ← (EXTI1) <-------/
            |---- gpio_in[2] ← (EXTI2)
            \---- ...
            \---- irq_out[7] → NVIC IRQ 7 (EXTI1_IRQn)
    \end{verbatim}
    \caption{qemu hierarchy}
    \label{fig:qemu_hierarchy}
\end{figure}

This hierarchy demonstrates the signal flow from the LIS3DH interrupt outputs through the EXTI (External Interrupt) controller to the NVIC (Nested Vectored Interrupt Controller), which ultimately triggers the ARM Cortex-M4 core's interrupt service routine.